\subsection{Lista de criterios}\label{subsec:dar-lista-criterios}

\subsubsection{Autenticación centralizada}

Este criterio determina si la tecnología analizada permite centralizar la autenticación de usuarios en una única plataforma, posibilitando el acceso controlado a múltiples servicios desde una sola fuente de identidad. Será evaluado como cumplimiento total (5) o inexistente (0).

\subsubsection{Autenticación multifactor (\texorpdfstring{\acs{MFA}}{MFA})}

Este criterio determina si la tecnología permite implementar mecanismos de autenticación multifactor, incorporando uno o más factores adicionales al uso de contraseña, como tokens, aplicaciones autenticadoras, certificados digitales o llaves físicas. Será evaluado según el nivel de soporte nativo o integrable, en una escala de 0 a 5.

\subsubsection{Ciclo de vida de cuentas}

Este criterio determina si la tecnología permite administrar el ciclo de vida de las cuentas de usuario, incluyendo creación, modificación, suspensión, expiración, reactivación y eliminación. Será evaluado de 0 a 5 según el nivel de automatización y control disponible.

\subsubsection{Cifrado de comunicaciones (\texorpdfstring{\acs{TLS}/\acs{LDAPS}}{TLS/LDAPS})}

Este criterio determina si la tecnología permite proteger las comunicaciones mediante protocolos seguros como \ac{TLS}, \ac{SSL} o \ac{LDAPS}, garantizando la confidencialidad e integridad en el intercambio de credenciales y datos. Será evaluado de 0 a 5.

\subsubsection{Compatibilidad con infraestructura existente}

Este criterio determina el nivel de compatibilidad de la tecnología con la infraestructura existente en el entorno analizado, particularmente con Ubuntu Server 24.04 \acs{LTS}, OpenVPN, Xen Orchestra y OpenProject. Será evaluado de 0 a 5 según la facilidad de integración y el soporte comprobado.

\subsubsection{Control granular de roles y permisos}

Este criterio determina si la solución permite asignar permisos diferenciados por usuario, grupo, rol, recurso o servicio, aplicando el principio de mínimo privilegio. Será evaluado de 0 a 5 según el nivel de detalle y flexibilidad disponible.

\subsubsection{Documentación y soporte técnico}

Este criterio determina la disponibilidad y calidad de documentación oficial, manuales, guías de instalación, foros comunitarios y soporte técnico empresarial o comunitario. Será evaluado de 0 a 5.

\subsubsection{Facilidad de administración}

Este criterio determina el nivel de simplicidad para instalar, configurar, operar, monitorear y mantener la plataforma, considerando herramientas gráficas, automatización y complejidad operativa. Será evaluado de 0 a 5.

\subsubsection{Gestión de usuarios y grupos}

Este criterio determina si la tecnología permite administrar usuarios y grupos de manera centralizada, facilitando su organización lógica, asignación de permisos y escalabilidad administrativa. Será evaluado de 0 a 5.

\subsubsection{Gestión y revisión de accesos}

Este criterio determina si la solución permite administrar autorizaciones y facilitar la revisión de los accesos para validar la vigencia de los privilegios, revocar permisos innecesarios y mantener el control institucional. Será evaluado de 0 a 5.

\subsubsection{Interoperabilidad entre sistemas}

Este criterio determina la capacidad de la tecnología para integrarse con diferentes aplicaciones, servicios y protocolos estándar como \ac{LDAP}, Kerberos, \ac{SAML}, \ac{RADIUS}, OpenID Connect y otros mecanismos de integración o federación. Será evaluado de 0 a 5.

\subsubsection{Políticas de contraseña}

Este criterio determina si la solución permite definir políticas de seguridad para contraseñas, tales como longitud mínima, complejidad, historial, expiración, bloqueo por intentos fallidos y mecanismos de recuperación segura. Será evaluado de 0 a 5.

\subsubsection{Registro y auditoría}

Este criterio determina si la tecnología permite registrar eventos relevantes como inicios de sesión, cambios administrativos, errores, intentos fallidos y modificaciones de permisos, facilitando la trazabilidad y auditoría. Será evaluado de 0 a 5.
